\paragraph{AI Alignment and Safety.} The field of AI alignment seeks to ensure that advanced AI systems pursue goals aligned with human values and interests.~\citet{bostrom2014superintelligence} articulated the control problem: a superintelligent agent, if mis-specified, could transform the world in undesirable ways. Such agents may exhibit instrumental convergence, that is, they will seek survival and resource acquisition as subgoals, unless explicitly prevented. Misalignment can arise because human values are complex and hard to specify formally, a problem also known as the \emph{outer alignment problem}, and because a system might develop its own emergent objectives --the \emph{inner alignment} or ``mesa-optimizer'' problem~\citep{hubinger2019risks}. Modern AI Alignment research has branched into several complmentary approaches. The closest alignment technique to the setup we study is Scalable Oversight~\citep{bowman2022measuring}. Scalable Oversight refers to the problem of designing oversight mechanisms that continue to work even as AI systems surpass human intelligence. The leveraged assumption is discrimination being easier than generation, thus spurring protocols such as debate~\citep{irving2018ai} and critiques~\citep{saunders2022self, mcaleese2024llm}. Debate in particular has been studied across many different forms and under different sets of assumptions. Our setup involves a single-turn sequential debate where the helpful and sneaky provers' goal is to persuade the verifier with their answer. Previous works on debate involve~\citet{michael2023debate}, which showed that debate outperforms consultancy, a protocol that involves showing the argument from a helpful prover half the time and sneaky half the time with human judges.~\citet{khan2024debating} showed that LLMs trained to maximize persuasiveness can lead to more truthful replies. Generally, empirical debate works have leveraged the QuALITY dataset~\citep{pang-etal-2022-quality}, which allows to model the asymmetry between debaters (or provers) and judges (or verifiers) as a function of access to information, that is, the debaters have unbounded access to data sources (long-form writeups), whereas the judges do not. Instead, we focus on formally verifiable domains such as code~\citep{hendrycks2021measuring}. The biggest source of inspiration for our work is~\citet{kirchner2024prover}, from which we borrow the Prover-Verifier setup: just like in our setup, they model the game as a single-turn interaction between two capable provers refereed by a weaker verifier.

\paragraph{AI Control.} Although traditional alignment research often assumes AIs to be at least trying to behave well after training, AI control poses a different question: \emph{can we remain safe even if the AI actively tries to subvert our safeguards?} In other words, instead of aligning the AI’s incentives, AI Control designs a second layer of defense that would catch or prevent catastrophes even in the event of the AI intentionally scheming. The practice of establishing safety measures that are robust to intentional subversion was first introduced in~\citet{greenblatt2023ai}. They introduced AI Control in the APPS backdoor setting, which our work builds on by relaxing the assumption of perfect strategy and perfect best response. An AI Control protocol is a formalization composed by the interplay of weak, trusted models, namely models that can be confidently claimed to be well aligned-or ``weak'' enough to not be able to reliably scheme-, and stronger, untrusted models, for which it is not assumed alignment.~\citet{griffin2024games} attempts a formalization, defining AI-Control Games: a kind of multi-agent decision process modeling the interaction between a blue team (the protocol designers) and a red team (a proxy for the scheming AI). In their formulation, the blue team’s policy (the deployment protocol) and the red team’s policy (an attack or subversion strategy) are engaged in a zero-sum partially observable stochastic game, with two objectives: safety (preventing catastrophic outcomes) and usefulness (getting useful work out of strong, untrusted models). By analyzing this as a game, one can formally search for equilibrium strategies or apply minimax optimization to find protocols that minimize worst-case risk while maximizing performance. Both~\citet{greenblatt2023ai} and~\citet{griffin2024games} consider optimal strategies to subvert control protocols but don't explore a model’s ability to generate the strategies themselves.~\citet{mallen2024subversion} investigate the capabilities an AI would need in
order to successfully subvert a control protocol statelessly.~\citet{mallen2024subversion} relaxed the assumption of perfect strategy, and delegates both the act of strategizing and execution to the AI. Importantly, they focus on a stateless setting: the AI is reset between tasks (no long-term memory), making covert planning harder.
\\
AI Control is more broadly inspired by adversarial robustness~\citep{goodfellow2019research}, and meta-level adversarial evaluations~\citep{shlegeris2023meta}, the practice of constructing adversarial policies for oversight processes. Adversarial robustness to a variety of inputs has been thoroughly studied in deep learning history~\citep{szegedy2013intriguing, biggio2013evasion, goodfellow2014explaining}, including recently in Large Language Models. Further work has explored automated monitoring systems to prevent jailbreaks of LLMs~\citep{zhang2023defending}, but these focus on adversarial users rather than the AI itself being adversarial. To the best of our knowledge, ours is the first work aimed at training subversion and oversight policies for the purpose of AI Control.

\paragraph{Computational complexity} Prover-Verifier Games borrow significantly from computational complexity, specifically interactive proof systems. In an interactive proof, a powerful, potentially unbounded \emph{prover} tries to convince a polynomial-time \emph{verifier} of a statement. Classic results like~\citet{babai1985trading} on Arthur-Merlin Games and~\citet{goldwasser1985ip} on interactive proofs showed that a clever protocol can extend what a limited verifier can confidently verify. Several complexity-theoretic concepts map onto alignment scenarios. A straightforward one is viewing an AI’s high-stakes decision or solution as a theorem that needs proof. Analogies can be made to NP vs. co-NP and to zero-knowledge proofs.~\citet{micali2000computationally} introduced computationally sound proofs, showing how a resource-bounded prover cannot cheat a verifier except with negligible probability. More concretely, recent work has applied interactive proof ideas to machine learning verification.~\citet{goldwasser2021interactive} and~\citet{pmlr-v195-mutreja23a} devised interactive protocols for Probabilistic-Approximately Correct (PAC) model verification, where a learning algorithm’s hypothesis can be checked via interaction. They show how a verifier can learn a hypothesis class with much less data by interacting with an untrusted prover. In both our work and~\citet{kirchner2024prover}, we assume capability, defined as size of model parametrization, (instead of data) constraint on the verifier. 
\\
The concept of AI Prover-Verifier Games is essentially an alignment-motivated re-imagination of interactive proofs. The first formalization of complexity theory protocols in this regard was made by Anil et al. (2021). The setup was a two-player game between a Prover (which produces an answer and explanation) and a Verifier (which approves or rejects the answer) By training both adversaries, they demonstrated the equilibrium of this game can favor truthful, checkable
answers: the verifier becomes robust to the sneaky prover’s tricks, and the truthful prover learns to present information in a way that even a limited verifier can validate. They also proved that different game formulations (e.g. who moves first, how the loss is structured) lead to different equilibria quality. Importantly, they found a Stackelberg equilibrium (leader-follower setup) where if the verifier commits to a strategy first, a rational prover’s best response is to be truthful, establishing theoretical conditions for \emph{completeness} and \emph{soundness}.~\citet{kirchner2024prover} is the closest work in terms of formalization to ours. They showed
that optimizing an AI solely for correctness can degrade the clarity of its reasoning, whereas optimizing for a verifier’s approval produces solutions that are easier for humans to check, albeit at a cost. They leverage the Prover-Verifier framework for training Language Models for legibility in the context of grade school math problems~\citep{cobbe2021training}, justifying the use of a weaker verifier as correlating with a human, who is bounded in capabilities with respect to a highly capable, possible superintelligent set of provers. Unlike in their formulation, our verifier’s task is not to score perceived correctness, but rather to classify the inputs as to whether they come from a honest distribution versus a malicious one.~\citet{amit2024models} trained LLM provers against a manually-defined verifier; in contrast, our verifier is learned on the basis of the aforementioned classification task. Notable is also the concept of Neural Interactive Proofs~\citep{hammond2024neural}, which generalizes PVGs to systems of neural networks. Their framework considered a bounded, trusted verifier (e.g., a smaller model or human simulation) interacting with one or more powerful, untrusted provers, all represented as neural networks. They unify various prior proposals (including debate, amplification, and~\citet{anil2021learning}'s PVG) under one umbrella and analyze them theoretically. By viewing these interactions through the formal lens of interactive proofs and zero-knowledge, they identify which protocols are more sample efficient or robust. They also introduce new interaction protocols and test them on tasks like graph isomorphism and code validation.

\paragraph{Adversarial Networks} Adversarial network methods provide a unifying framework for control-oriented training by pitting ``red'' (attacker or subversion) models against ``blue'' (defender or oversight) models in a minimax game. The canonical example is the Generative Adversarial Network (GAN) of \cite{goodfellow2014generative}, which trains a generator $G$ to fool a discriminator $D$ in a two-player adversarial process.  This basic idea has been adapted in many safety‐critical settings: for instance,~\citet{pinto2017robust} formulate robust reinforcement learning (RARL) as a two‐agent zero‐sum game where a control policy (the ``protagonist'') learns to perform a task while a learned adversary agent simultaneously learns to perturb the environment or apply disturbances to make the task harder, yielding policies that are empirically more stable and robust to modeling errors. Likewise, standard adversarial‐training in supervised learning (e.g., FGSM‐based methods) explicitly augments training data with worst‐case perturbations to harden models;~\citet{goodfellow2014explaining} show that incorporating adversarial examples in training can significantly improve model robustness against such attacks. In parallel,~\citet{ho2016generative} repurposes policy learning from demonstrations as a GAN‐like process: a generator policy attempts to produce expert‐like behavior that a discriminator cannot distinguish from real demonstrations, effectively learning a policy that both imitates and is verified by the discriminator.
\\
In AI safety research, adversarial techniques have been explicitly exploited for oversight and worst‐case analysis. For example,~\citet{ziegler2022adversarial} show that adversarial training, including tools‐assisted human adversaries, can greatly boost the reliability of a language-filter classifier by systematically discovering and removing dangerous outputs. More abstractly,~\citet{kirchner2024prover} can be thought of as a form of adversarial game. This setup forces the verifier to become robust to sophisticated attacks, thereby improving the provers’ checkability (legibility) and uncovering subtle failure strategies.

\paragraph{Game Theory} Game theory provides a natural lens for many AI alignment and control problems, which often involve multiple agents (human, AI, or subsystems of AI) with partially or fully opposing goals. A straightforward application is viewing AI oversight as a zero-sum game between a would-be misaligned AI and an oversight mechanism (human or AI). The oversight wants to minimize bad behavior; the misbehaving policy, in the worst case, wants to maximize it. Formulating this as a minimax optimization aligns with robust control: we train the overseer (or the agent’s policy) to handle the worst possible strategy of an adversary. This is evident in adversarial training setups and was made explicit in works like~\citet{pinto2017robust}, which amounts to solving a minimax game.
\\
Another aspect is mechanism design and incentives. From a game theory viewpoint, if we treat an AI as an agent with its own utility, the alignment problem becomes one of incentive design: how do we structure the ``game'' (be it the training process, the reward function, the environment) such that the AI’s best move is to act in alignment with human values? Prover-Verifier Games themselves are essentially a two-player game framework. The prover and the verifier can be seen as players with opposite objectives (one tries to convince, the other to only be convinced of
true statements). In Stackelberg PVG~\citep{anil2021learning}, the verifier commits to a strategy (like a rule for accepting answers), and the prover then best-responds. If the verifier’s commitment is optimal, it will be one that makes truthful proving the prover’s best option.~\citet{anil2021learning} showed that under reasonable assumptions there is a desirable equilibrium (honest behavior) in their game, and empirically they converged near it. Furthermore, iterated games where in each round the red team (trying to break the AI or protocol)
finds a weakness, then the blue team adjusts the AI or protocol to patch that, are analogous to concepts like an arms race or an iterative minimax. In fact,~\citet{griffin2024games} treats the red-team optimization and blue-team optimization as intertwined, and search for an equilibrium of that combined game. Solving (or approximately solving) these games can tell us, for instance, whether a given protocol is likely to succeed against a strategic attacker, or whether the attacker has a winning strategy. It can also quantify the cost of safety: e.g., how much utility is lost at equilibrium due to the constraints needed to block attacks (sometimes called a price of security).